\documentclass[11pt]{article}
\usepackage[margin=1in]{geometry}
\usepackage{amsmath,amssymb,amsthm}
\usepackage{xurl}
\sloppy
\let\oldus\_ \renewcommand{\_}{\oldus\allowbreak}
\newtheorem{theorem}{Theorem}
\newtheorem{lemma}[theorem]{Lemma}
\theoremstyle{definition}
\newtheorem{definition}[theorem]{Definition}
\title{Conjecture 00000000539 is false:\\ the h-vector $(1,4,6,4,1)$ of $(x_0^2,x_1^2,x_2^2,x_3^2)$}
\author{Jackmeson1}
\date{}
\begin{document}
\maketitle

\section{The conjecture}
\emph{Definition: The h-polynomial $h(t)$ is the encoding generating function of the Hilbert
function of a graded ideal. Conjecture: The h-vector of a quadratic monomial ideal has no three
adjacent strictly decreasing terms; this combinatorial positivity is induced by the shellability
of the complex.}

The conjecture is a conjunction. Its first clause is a universal statement about all quadratic
monomial ideals. We refute that clause, so the conjunction is false. The second clause (the
"explanation" by shellability) is not separately formalized.

\section{Reading and definitions}
Let $K$ be any field and $S=K[x_0,\dots,x_{n-1}]$ with the standard grading ($\deg x_i=1$).
For a set $G$ of exponent vectors, the \emph{monomial ideal} generated by $G$ is
$I_G=(x^g : g\in G)$. A \emph{quadratic monomial ideal} is an ideal $I_G$ with every $g\in G$ of
total degree $2$. Squares $x_i^2$ are allowed; nothing in the text requires squarefree generators.

\begin{definition}[Hilbert function, Hilbert series, h-polynomial]
For a homogeneous ideal $I$, $(S/I)_j$ is the image of the degree-$j$ homogeneous polynomials
$S_j$ under $S\to S/I$, and $\mathrm{HF}_{S/I}(j)=\dim_K (S/I)_j$. The Hilbert series is
$\mathrm{HS}_{S/I}(t)=\sum_j \mathrm{HF}_{S/I}(j)\,t^j\in\mathbb{Z}[[t]]$. A polynomial
$h\in\mathbb{Z}[t]$ is \emph{the h-polynomial} of $S/I$ if
$\mathrm{HS}_{S/I}(t)\,(1-t)^d=h(t)$ for some $d\ge 0$ and $h(1)\neq 0$. Its coefficient list
$(h_0,\dots,h_s)$, $s=\deg h$, is the h-vector.
\end{definition}

\noindent\textbf{Conventions.} (1) Reading A (standard): ``the Hilbert function of a graded
ideal $I$'' is that of the quotient $S/I$. (2) Reading B (literal): it is the Hilbert function
$\mathrm{HF}_I(j)=\dim_K I_j$, $I_j=I\cap S_j$, of the graded module $I$, with h-polynomial
defined by the same rule from $\mathrm{HS}_I$. (3) ``Three adjacent strictly decreasing terms''
means an index $i$ with $i+2\le\deg h$ and $h_i>h_{i+1}>h_{i+2}$. Coefficients are integers.
(4) $\mathbb{N}$ contains $0$.

\section{The counterexample}
Let $n=4$ and $I=(x_0^2,x_1^2,x_2^2,x_3^2)$, a quadratic monomial ideal.

\begin{lemma}[standard monomials]\label{lem:std}
For a monomial ideal $I_G$ in finitely many variables, $\dim_K (S/I_G)_j$ equals the number of
exponent vectors $m$ of degree $j$ such that no $g\in G$ satisfies $g\le m$ coordinatewise.
\end{lemma}
\begin{proof}
A polynomial lies in $I_G$ if and only if every monomial in its support is divisible by some
$x^g$, $g\in G$ (Mathlib: \texttt{MvPolynomial.mem\_ideal\_span\_monomial\_image}). Hence a
nonzero combination of standard monomials (those divisible by no $x^g$) is never in $I_G$, so
their images in $S/I_G$ are linearly independent. Every non-standard monomial maps to $0$, and
$S_j$ is spanned by the degree-$j$ monomials. So the images of the degree-$j$ standard monomials
form a basis of $(S/I_G)_j$.
\end{proof}

\begin{theorem}[Reading A]
$\mathrm{HF}_{S/I}(j)=\binom{4}{j}$, so $\mathrm{HS}_{S/I}(t)=(1+t)^4$. The h-polynomial is
uniquely $(1+t)^4$, with h-vector $(1,4,6,4,1)$, and $6>4>1$ are three adjacent strictly
decreasing terms.
\end{theorem}
\begin{proof}
For $G=\{2e_0,\dots,2e_3\}$, $m$ is standard iff every $m_i\le 1$. Such $m$ of degree $j$
correspond bijectively to their supports, the $j$-subsets of $\{0,1,2,3\}$, so
Lemma~\ref{lem:std} gives $\binom4j$. Then $(1+t)^4\cdot(1-t)^0$ is a valid h-polynomial,
$h(1)=16\ne0$. If $\mathrm{HS}(t)(1-t)^d=h$ with $h(1)\neq0$, then $h=(1+t)^4(1-t)^d$ as
polynomials and $h(1)=16\cdot 0^d$, so $d=0$ and $h=(1+t)^4$.
\end{proof}

\begin{theorem}[Reading B]
$\mathrm{HF}_I(j)=\binom{j+3}{3}-\binom4j$ and $\mathrm{HS}_I(t)(1-t)^4=1-(1+t)^4(1-t)^4
=4t^2-6t^4+4t^6-t^8=:h_I(t)$, with $h_I(1)=1\ne 0$. Every h-polynomial of $I$ equals $h_I$; its
h-vector $(0,0,4,0,-6,0,4,0,-1)$ has $4>0>-6$ in positions $2,3,4$.
\end{theorem}
\begin{proof}
Rank--nullity for $S_j\to(S/I)_j$ (kernel $I_j$) gives $\dim I_j+\dim(S/I)_j=\dim S_j$, and
$\dim S_j=\binom{j+3}{3}$ (number of degree-$j$ exponent vectors in four variables). With
$\sum_j\binom{j+3}{3}t^j\cdot(1-t)^4=1$ the identity follows. Uniqueness: if
$F(1-t)^a=p$ and $F(1-t)^b=q$ with $a\le b$ and $q(1)\neq0$, then $q=p(1-t)^{b-a}$, forcing
$a=b$ and $p=q$.
\end{proof}

\noindent\textbf{Remark (not formalized).} The squarefree polarization
$J=(x_0y_0,x_1y_1,x_2y_2,x_3y_3)\subset K[x_i,y_i]$ (edge ideal of a perfect matching on 8
vertices; its Stanley--Reisner complex is the boundary of the 4-dimensional cross-polytope)
has the same h-vector $(1,4,6,4,1)$: its degree-$j$ standard monomials correspond to pairs
(subset $s$ of blocks with a $y$-power, arbitrary exponent vector $b$ of degree $j-|s|$), so
$\mathrm{HS}_{S/J}=(1+t)^4/(1-t)^4$. This is only a paper argument; the Lean development covers
the ideal $I$ above.

\section{Lean formalization}
File \texttt{lean/Conjecture539/Basic.lean} (Lean 4, Mathlib v4.33.1, only \texttt{import Mathlib}).
\begin{itemize}
\item Objects: \texttt{monomialIdeal}, \texttt{IsQuadraticMonomialIdeal}, \texttt{gradedPiece}
 (image of \texttt{homogeneousSubmodule} in $S/I$), \texttt{hilbertFunction} (\texttt{finrank}),
 \texttt{hilbertSeries}, \texttt{IsHPolynomial}, \texttt{HasTripleDescent}; for reading B
 \texttt{idealPiece}, \texttt{idealHilbertFunction}, \texttt{idealHilbertSeries},
 \texttt{IsIdealHPolynomial}.
\item \texttt{hilbertFunction\_monomialIdeal}: Lemma~\ref{lem:std} for every finite variable set.
\item \texttt{hilbertFunction\_I4}, \texttt{hilbertSeries\_I4}, \texttt{isHPolynomial\_I4},
 \texttt{isHPolynomial\_unique}, \texttt{h4\_tripleDescent}, \texttt{counterexample}.
\item \texttt{idealHF\_add\_HF}, \texttt{finrank\_homogeneous\_four},
 \texttt{idealHilbertSeries\_I4\_mul}, \texttt{isIdealHPolynomial\_I4},
 \texttt{isIdealHPolynomial\_unique}, \texttt{hI\_tripleDescent}.
\item Main theorems, for every field $K$:\\
\texttt{conjecture\_false (K) : Not (forall n (I : Ideal (MvPolynomial (Fin n) K)), IsQuadraticMonomialIdeal I -> forall h, IsHPolynomial I h -> Not (HasTripleDescent h))}\\
and \texttt{conjecture\_false\_ideal\_reading}, the same with \texttt{IsIdealHPolynomial}.
\item Axiom audit: \texttt{\#print axioms} for both main theorems and \texttt{counterexample}
 reports only \texttt{propext}, \texttt{Classical.choice}, \texttt{Quot.sound}. No \texttt{sorry},
 no \texttt{native\_decide}.
\end{itemize}

\end{document}
